\chapter{绪论}

\section{研究背景与意义}
脑机接口（Brain-Computer Interface, BCI）是在大脑与外部设备之间建立信息交换通路的技术，近年来正从实验室验证逐步进入医疗康复、人机协同、智能控制和娱乐交互等应用场景\cite{li2025bciindustry,meng2025adversarialbci}。按照信号采集方式划分，BCI 可分为侵入式、半侵入式和非侵入式三类。非侵入式 BCI 以脑电信号（Electroencephalogram, EEG）为主要输入，具有部署成本低、使用门槛低和安全性较高等特点，适合与可穿戴设备、移动终端和边缘节点结合。

物联网（Internet of Things, IoT）为 BCI 系统提供了联网部署基础。典型 BCI-IoT 系统由采集端、无线接入链路、边缘网关、云端服务和反馈执行端构成。采集端持续上传脑电片段、分类结果或控制请求，服务端完成信号处理、状态分析和反馈决策，再将反馈指令返回终端或执行设备。该模式扩展了 BCI 的服务范围，也使神经信号和控制指令长期暴露在无线网络、公网链路、代理节点和云端服务组成的复杂通信路径中。

BCI-IoT 的通信安全具有鲜明的场景特征。首先，脑电信号与用户身份、认知状态、任务类别和控制意图存在关联，泄露后难以像普通口令一样更换。其次，采集端多为可穿戴或边缘设备，计算能力、内存容量和能耗预算受限，安全协议需要控制握手开销和运行时状态规模。再次，BCI 闭环控制依赖上行信号和下行反馈的连续交互，重放、篡改和中间人攻击可能影响控制结果\cite{meng2025adversarialbci,liu2022accesscontrol,trabelsi2023iotaccess}。因此，面向 BCI-IoT 的安全通信协议需要同时考虑会话安全、资源开销、可部署性和流量外观暴露。

传统安全通道可以保护业务明文，但包长、方向、到达间隔和字节统计等外部特征仍可被长期观测并用于流量侧写\cite{velan2015survey,alwhbi2024encrypted,hou2024encmalware,fu2025encryptedtraffic}。在 BCI-IoT 场景中，不同脑电任务、反馈频率和控制状态可能对应不同的消息长度和交互节奏，进而形成可学习的外观指纹。本文围绕 BCI-IoT 链路中的身份认证、会话保护、抗重放、异常处置和外观抗侧写需求，设计并实现一种基于组合数独外观编码的安全通信协议原型，为受限设备上的 BCI-IoT 安全通信提供可复现的工程验证。

\section{国内外研究现状}
\subsection{脑机接口安全与隐私研究}
BCI 安全研究已覆盖隐私泄露、对抗样本、系统攻击、伦理治理和安全框架等多个方向。Lopez-Bernal 等对 BCI 生命周期中的采集、处理、分类和反馈阶段进行了攻击面归纳，指出 BCI 系统的机密性、完整性、可用性与人身安全存在紧密联系\cite{lopezbernal2021bcisecurity}。Martinovic 等验证了消费级 BCI 设备中脑电信号被用于隐私推断的可行性\cite{martinovic2012sidechannel}，Lange 等进一步以 PIN 输入场景分析了脑信号侧信道攻击的实验条件和成功因素\cite{lange2018pinbci}。这些研究表明，脑电数据、任务类别和控制意图均可能成为攻击者关注的对象。

随着 BCI 设备接入物联网，安全问题从单一终端扩展到端、边、云协同系统。Tarkhani 等对 Muse、NeuroSky、OpenBCI 等可穿戴 BCI 产品进行分析，报告了硬件、软件和网络栈中的多类漏洞，并提出信息流控制机制以限制脑电数据泄露\cite{tarkhani2022wearablebci}。Ajmeria 等从工业物联网视角综述了基于 EEG 的 BCI 应用，指出 BCI-IoT 可将人类认知状态引入工业控制和智能系统\cite{ajmeria2023iiotbci}。现有研究为 BCI-IoT 安全奠定了风险认知基础，但针对通信协议层的细粒度实证研究仍相对有限。

\subsection{物联网安全通信协议研究}
物联网安全通信常采用 TLS、DTLS、MQTT/TLS、CoAP/DTLS 等标准化方案。TLS 1.3 面向 TCP 通信，DTLS 1.3 面向 UDP 数据报通信，二者通过身份认证、密钥协商和记录层保护为业务数据提供安全通道\cite{rfc8446,rfc9147}。相关研究指出，在低功耗设备和受限网络中部署 TLS/DTLS 时，需要关注握手轮次、证书管理、内存占用、UDP 可达性和链路重传等工程因素\cite{restuccia2020lowpower,nedergaard2023coapsecure}。

MQTT 与 CoAP 是物联网系统中常见的应用层协议。MQTT 基于发布/订阅模式，适用于低带宽和弱连接环境，通常借助 TLS 建立安全通道\cite{mqtt311,zheng2023mqtttamarin}。CoAP 面向受限节点与受限网络，运行在 UDP 之上，并通过消息 ID、Token 与 CON/ACK 机制支持轻量请求响应交互\cite{rfc7252,ren2020coapcongestion}；在需要加密认证时，CoAP 通常与 DTLS 结合使用\cite{rfc9147,nedergaard2023coapsecure}。上述协议生态成熟，适合作为 BCI-IoT 安全通信的工程参照。

除标准协议组合外，部分轻量系统会采用自定义记录层或仅使用认证加密（Authenticated Encryption with Associated Data, AEAD）保护数据\cite{rfc5116}。该类方案可降低实现复杂度和运行开销，但身份认证、密钥新鲜性、重放检测和异常终止条件需要由上层状态机补足。对于 BCI-IoT 场景，单纯降低开销不足以覆盖脑电数据敏感性和闭环控制风险，协议设计还需要考虑攻击输入、终止代理、探测流量和外观侧写带来的复合威胁。

\subsection{加密流量侧写研究}
加密流量分类研究表明，即使载荷内容已经加密，长度、时序、方向和字节分布等外部特征仍可用于协议识别、恶意流量检测和用户行为推断\cite{velan2015survey,alwhbi2024encrypted,hou2024encmalware,fu2025encryptedtraffic}。在 BCI-IoT 链路中，脑电任务类别、反馈频率和控制状态会影响消息长度、发送间隔和上下行比例，攻击者可将这类特征作为侧写输入，从而推断用户状态或控制意图。

现有抗侧写方法主要包括固定长度记录、随机填充、流量整形、批量发送和时序扰动等。这些方法能够削弱外观特征与业务类别之间的稳定关联，但会引入额外带宽、时延和能耗成本。面向资源受限 BCI-IoT 终端，外观防护需要与会话安全、端侧资源和实际链路可达性共同评估。本文将随机填充与可逆组合数独外观编码引入协议数据面，并通过协议识别、类别恢复和 DPI 规则实验评估其效果。

\section{现有问题分析}
结合上述研究，BCI-IoT 安全通信仍存在以下问题。

第一，BCI 安全研究对神经数据隐私、模型攻击和伦理治理已有较多讨论，但通信协议层的工程验证相对不足。联网 BCI 系统中的攻击者可同时观察网络外观、干扰握手流程和注入探测流量，单纯讨论数据加密难以覆盖闭环控制链路中的重放、篡改和异常输入风险。

第二，物联网安全协议研究通常以通用传感器数据或设备遥测为对象，关注握手时延、内存占用和吞吐能力。BCI-IoT 传输的内容具有更强的个人敏感性和控制语义，长度序列、发送节奏和字节分布可能暴露脑电任务类别。现有标准通道能够保护明文，却较少将外观侧写能力纳入同一评测框架。

第三，轻量化记录层方案有利于资源受限设备部署，但仅依赖 AEAD 记录保护时，认证握手、密钥新鲜性、重放缓存和异常处置仍需额外设计。BCI-IoT 终端一旦接入公网或跨域网络，还会遇到 UDP 受限、代理终止、策略审查和探测流量等部署问题。协议设计需要给出可执行状态机，并通过实验说明安全收益与工程代价。

基于上述问题，本文选择在应用层与 TCP 之间实现传输增强层。该层以认证握手建立会话身份与密钥，以 AEAD 记录层保护业务明文，以组合数独外观层调节可观测字节分布和长度特征，并在本地回环、侧写攻防和跨区域受限网络中进行对照验证。

\section{研究内容与创新点}
本文围绕 BCI-IoT 安全通信需求开展协议设计、原型实现和实验评估，主要研究内容如下：
\begin{itemize}
  \item 设计面向 BCI-IoT 的通信层安全模型，明确脑电数据、控制指令、会话材料、协议状态和外观特征对应的安全属性；
  \item 实现认证握手与 AEAD 记录层，完成证书签名、临时密钥协商、摘要绑定、确认码、重放检测和异常路径处置；
  \item 设计可逆组合数独外观层，将业务字节映射为提示字节序列，并结合随机填充调节可打印比例、字节熵和长度序列；
  \item 搭建统一实验工具链，与 DTLS、MQTT/TLS、CoAP/DTLS 和纯 AEAD 方案进行性能、资源、侧写和主动攻击对照；
  \item 在本地回环、受控侧写实验和跨区域受限网络中验证协议原型的可用性与适用条件。
\end{itemize}

本文的主要创新点包括：
\begin{itemize}
  \item 面向 BCI-IoT 场景，将会话安全、资源约束和外观抗侧写放入同一协议原型中验证，补充了通信层实证研究；
  \item 提出基于组合数独的可逆外观编码方法，在保持解码确定性的同时改变可打印比例、字节熵和长度序列特征；
  \item 构建包含重放、篡改、探测流量、协议侧写、BCI 类别恢复和 TLS 中间人对照的实验流程，给出协议收益与代价的量化结果。
\end{itemize}

\section{论文结构安排}
全文共五章，结构如下：
\begin{itemize}
  \item 第 1 章为绪论，介绍研究背景与意义、国内外研究现状、现有问题、研究内容和论文结构；
  \item 第 2 章为定义、模型与理论基础，说明 BCI-IoT 通信对象、安全资产、威胁模型、研究范围、基础密码机制、流量侧写指标和组合数独外观层数学基础；
  \item 第 3 章为协议设计与实现，给出认证握手、记录层、组合数独外观层、异常路径和工程实现；
  \item 第 4 章为实验设计与结果分析，说明实验环境、指标计算、性能结果、侧写结果、主动攻击验证和工程权衡；
  \item 第 5 章为总结与展望，总结本文工作并讨论后续研究方向。
\end{itemize}
